\documentclass{amsart}
\newtheorem{theorem}{Theorem}[section]
\newtheorem{lemma}[theorem]{Lemma}
\newtheorem{definition}[theorem]{Definition}
\numberwithin{equation}{section}
\newcommand{\dbl}{\mathrm{d}}
\begin{document}
\section{Even numbers}\label{sec:even}
Throughout, $m$, $n$ and $k$ are natural numbers.
\begin{definition}\label{def:even}
A number $n$ is \emph{even} if $n=2k$ for some $k$. We write $\dbl(k)=2k$.
\end{definition}
\begin{lemma}\label{lem:double}
For every $k$, the number $\dbl(k)$ is even.
\end{lemma}
\begin{proof}
By definition, $\dbl(k)=2k$ is of the required form.
\end{proof}
\begin{lemma}\label{lem:sum}
If $m$ is even and $n=\dbl(b)$ for some $b$, then $m+n$ is even.
\end{lemma}
\begin{proof}
Write $m=2a$. By Lemma~\ref{lem:double}, $n$ is even, i.e.\ $n=2b$.
Then
\begin{equation}\label{eq:sum}
m+n=2a+2b=2(a+b).
\end{equation}
This is the familiar argument.
\end{proof}
\end{document}
